HealthGraphBench is an exploratory benchmark devoted to three temporal health-regulatory tasks: new product--problem links in MAUDE, G--L deficiencies in CMS-monitored nursing-home inspections, and new prescriber--drug links in Medicare Part D. We distinguish historical benchmark results from newly tuned comparisons and earlier diagnostic analyses on MAUDE.

In the historical MAUDE results, neighbor frequency achieves micro Recall@10 of 0.192045, compared with 0.172840 for three-epoch GraphSAGE. The historical duration analysis C selects thirty epochs on validation and reaches 0.210642 on 2024--2025, compared with 0.192045 for neighbors; the difference of $+0.018597$ has a paired conditional 95\,\% CI of $[+0.011975;+0.025570]$. This CI does not transfer to the new Q2 fits. In this validation-tuned comparison, selected BPR achieves mean micro Recall@10 of 0.230551 in 2024 and 0.239437 in 2025, ahead of the fixed baselines. Mean-aggregation GraphSAGE reaches 0.204065 and 0.211825, respectively; the control without message passing reaches 0.193403 and 0.165388. Q2 values for stochastic models are averaged over the three seeds, with no seed selection. The historical D1 diagnostic retains a nearly flat loss; the Q2 results do not establish that this near-flatness resulted from a historical optimization bug.

On CMS, in the historical analysis, adding owner aggregates increased pooled AUC by 0.001318 (historical CI includes zero); Q2 comparisons show small differences that are not uniformly positive across algorithms and years. On Part D, Q2 BPR reaches 0.236182 on the 2024 test, compared with 0.217860 for specialty popularity; historically, this heuristic outperformed BPR (0.217860 versus 0.180789). Tunable families were selected on validation; the results remain exploratory because these periods had already been examined. No new interval is calculated for the Q2 fits. A separate prospective forecast concerns the future publication of Part D relationships, not a clinical outcome; independent evaluation has not yet been conducted. These results demonstrate neither general graph superiority nor clinical utility.
